\chapter{Campagne d'endurance du firmware}
\label{ann:endurance}

Cette annexe précise la méthode de la campagne d'endurance dont les résultats sont exploités à la sous-section~\ref{sec:archi_validation} et donne la répartition détaillée de la charge processeur.

\section*{Méthode de relevé}

Une tâche de diagnostic de faible priorité imprime toutes les dix secondes, sur la liaison série, un bloc récapitulant l'état du système. Ce bloc contient l'occupation des piles, obtenue pour chaque tâche par \texttt{uxTaskGetStackHighWaterMark} qui donne le plus petit espace de pile resté libre depuis le démarrage, la répartition du temps processeur fournie par les statistiques d'exécution de FreeRTOS, ainsi que le tas interne et la PSRAM encore disponibles. Les pourcentages de charge sont calculés par différence entre deux relevés consécutifs, la somme des charges de l'ensemble des tâches sur un même intervalle couvrant le temps des deux cœurs. La campagne complète compte 1259 relevés sur 3~h~34.
\section*{Répartition de la charge par régime}

Le tableau~\ref{tab:charge_regime} donne la charge moyenne de chaque tâche significative et de chaque cœur pour les régimes rencontrés au cours de la campagne. Il fait apparaître deux résultats. La charge de la tâche audio croît avec l'effort de décodage, du simple transfert d'un flux WAV au décodage complet d'un MP3. La pile Bluetooth n'apparaît que sur le cœur~0 pendant le streaming et n'ajoute qu'une charge marginale à la tâche audio, ce qui confirme l'isolement des deux cœurs. Les valeurs sont des moyennes par régime et sont donc inférieures aux crêtes instantanées visibles sur la figure~\ref{charge_cpu}.

\begin{table}[H]
    \centering
    \caption{Charge processeur moyenne par régime, en pour-cent du temps de chaque cœur. La charge du cœur inclut, en plus des tâches détaillées, les tâches système résiduelles restant sous 1~\%. La charge de la tâche audio suit le format décodé et la pile Bluetooth reste confinée au cœur~0.}
    \label{tab:charge_regime}
    \begin{tabular}{lccccc}
        \toprule
                                  & Repos & WAV 16 bits & WAV 24 bits & MP3  & Bluetooth \\
        \midrule
        \multicolumn{6}{l}{\textit{Cœur 1}}                                              \\
        \quad audio [\%]          & 0,0   & 11,5        & 16,1        & 36,8 & 38,8      \\
        \quad ui [\%]             & 2,0   & 0,4         & 0,3         & 0,7  & 1,2       \\
        \quad charge du cœur [\%] & 2,0   & 11,9        & 16,4        & 37,5 & 40,0      \\
        \midrule
        \multicolumn{6}{l}{\textit{Cœur 0}}                                              \\
        \quad maintenance [\%]    & 2,1   & 2,1         & 2,1         & 2,1  & 2,2       \\
        \quad pile Bluetooth [\%] & 0,0   & 0,0         & 0,0         & 0,0  & 67,0      \\
        \quad charge du cœur [\%] & 2,1   & 2,1         & 2,1         & 2,1  & 70,0      \\
        \bottomrule
    \end{tabular}
\end{table}


\section*{Stabilité du tas interne}

La figure~\ref{tas_interne} donne l'évolution du tas interne libre sur toute la campagne. Le niveau varie avec les allocations de la pile Bluetooth et des tampons associés mais ne présente aucune décroissance monotone, ce qui confirme l'absence de fuite mémoire. Les relevés échantillonnés toutes les dix secondes ne descendent pas sous 69~kio, tandis que le suivi continu du firmware rapporte un plancher de 55,6~kio en tenant compte des creux transitoires entre deux relevés.

\fig[H, width=1\textwidth]{Tas interne libre au cours de la campagne d'endurance. Le fond reprend le code couleur des régimes de la figure~\ref{charge_cpu}. Le niveau ne décroît pas dans la durée, ce qui écarte toute fuite mémoire.}{tas_interne}